% !TeX root = ../本科微分几何笔记.tex
% 由主文档通过 \input 载入；请编译主文档。
\chapter{绪论：研究对象、课程定位与学习路线}
\label{chap:dg-intro}

\section{从弯曲现象到几何不变量}
一条曲线可以用不同速度描画，一张曲面也可以用不同坐标描述。位置函数及其导数会随这些选择改变，但曲线的长度、曲面上两条切向量的夹角等几何量不应因此改变。\emph{微分几何（differential geometry）}用微积分和线性代数研究光滑几何对象，并把不依赖描述方式的性质组织成理论。本书研究欧氏空间中的曲线与曲面，逐步区分三个层次：参数表示是计算工具，切空间与度量是几何结构，曲率则刻画这些结构的变化。

纸片弯成圆柱以后，纸面上曲线的长度可以保持不变，而纸片在空间中的法向方向发生了变化。这个例子说明，\emph{内蕴几何（intrinsic geometry）}与\emph{外蕴几何（extrinsic geometry）}提出的问题不同：前者只依赖对象内部的长度与角度，后者还研究对象如何放置在周围空间中。曲面的第一基本形式描述前一种信息，第二基本形式描述后一种信息；高斯曲率同时出现在两种理论中，这一事实是全书的中心线索。

局部研究关注一个点附近的几何，整体研究关注整个对象。曲线各处的切向量转动，最终可以决定闭曲线的旋转指标；曲面的高斯曲率积分，则与欧拉示性数联系起来。局部公式不自动包含整体结论：圆柱局部可以展开成平面，却不能由一个单值、无接缝的平面坐标覆盖整个圆柱；测地线在足够短的区间内最短，却可能在走过较长距离以后失去最短性。

\section{参考体系与编写原则}
本书以北京大学公开的本科微分几何课程说明作为内容范围的主要参照，其先修课为数学分析、高等代数与常微分方程，教学主线由曲线和曲面的基本形式推进到曲面基本方程、测地线及 Gauss--Bonnet 公式\cite{pku-course}。在论证方式上，参考 MIT 18.950 的公开课程定位：以曲率为中心，结合具体几何与定义、定理、证明，并贯通局部几何、整体几何及长度距离问题\cite{mit-syllabus,mit-notes}。这里具体指 MIT 2008 年秋季公开课程，不把某一年的课程安排视作所有年份的统一要求。

清华大学数学科学系的公开教师课程信息也将本科微分几何与研究生微分流形课程分列\cite{tsinghua-course}；该信息用于确认课程层次，不据此推定其指定教材或完整教学大纲。因此，本书选择适合中文本科教学的曲线曲面路线，将抽象流形和黎曼度量放在末章衔接。陈维桓的《微分几何初步》、do Carmo 的曲线曲面教材和 K\"uhnel 的教材列为延伸阅读；这些书的书目信息可分别从北大和 MIT 的课程页面核验\cite{chen,docarmo,kuhnel}。

正文、例题与习题按知识依赖独立组织，不是上述教材的逐章翻译，也不代表有关学校的正式出版教材。基础定理尽量给出完整证明；需要额外拓扑或偏微分方程工具的结论，明确标记所引用的前置定理以及证明边界。判断一段论证是否成立，应检查假设、使用的定理与推理过程，而不是以参考学校的名称代替数学依据。

\section{先修知识与阅读顺序}
读者应熟悉一元与多元微积分、链式法则、逆函数与隐函数定理、二重积分及 Green 公式，掌握内积空间、对称矩阵与正定二次型，并了解常微分方程初值问题的局部存在唯一性。第~\ref{chap:dg-tools} 章集中说明反复使用的工具，不能替代完整的分析与代数课程。

第~\ref{chap:dg-curves}、\ref{chap:dg-plane} 章研究曲线的局部与整体性质；第~\ref{chap:dg-surfaces}--\ref{chap:dg-examples} 章建立曲面、度量、形算子及典型模型；第~\ref{chap:dg-connection} 章把曲面基本方程与内蕴曲率联系起来；第~\ref{chap:dg-geodesics}、\ref{chap:dg-gb} 章处理测地线与整体曲率；第~\ref{chap:dg-manifolds} 章给出后续流形课程的入口。等周不等式、曲面论基本定理和流形一章可作为提高内容，但不应跳过其假设说明。

每章设置四道自编习题，分别训练概念辨析、计算、证明与综合应用，解答集中于附录~\ref{app:dg-solutions}。学习时宜先独立解题，再对照解答检查论证；公式汇总放在附录~\ref{app:dg-reference}，只用于查阅，不能替代推导。

\section{正则性、符号与方向约定}
除另行说明，映射、曲线、曲面及向量场均假设光滑，即属于 $C^\infty$ 类；特定命题若说明只需 $C^k$，则按其较弱假设理解。这样可以在主线推导中合法交换混合偏导数，但不省略正则性、非退化性、紧致性与连通性等独立条件。曲线参数区间默认连通；讨论闭曲线时要求首尾各阶导数按所需阶数匹配。

欧氏内积记为 $\ip{v}{w}$，范数记为 $\norm{v}$；曲线用 $\gamma$ 表示，弧长参数用 $s$，一般参数用 $t$。曲面记为 $\Sigma$，局部参数映射记为 $\mathbf x(u,v)$，选定的单位法向量记为 $\nu$。形算子统一取 $S=-\dd\nu$，第二基本形式取 $\mathrm{II}(X,Y)=\ip{SX}{Y}$。因此，半径为 $R$ 的球面取外法向时主曲率为 $-1/R$；高斯曲率为 $1/R^2$。其他教材可能使用相反的形算子符号，比较公式前必须先核对约定。

度量系数为 $g_{ij}$，逆矩阵系数为 $g^{ij}$，两类指标不能混淆。涉及重复指标的求和通常显式写出；只有在明确声明后才采用 Einstein 求和约定。曲面定向、边界正向、曲线行进方向及空间刚体运动是否保持定向，均在相关结论中说明。

